\documentclass[12pt]{extreport}
\usepackage{indentfirst}
\usepackage[T2A]{fontenc}
\usepackage[utf8]{inputenc}
\usepackage[english,russian]{babel}
\usepackage{parskip}
\usepackage[top = 20mm, left = 30mm, right = 10mm, bottom = 20mm]{geometry}
\setlength\parindent{12.5mm}
\renewcommand{\baselinestretch}{1.5}
\def\hrf#1{\hbox to#1{\hrulefill}}
\usepackage{ragged2e}
\usepackage{tikz}
\usetikzlibrary{arrows.meta, positioning, shapes.geometric}
\usepackage{multicol}
\usepackage{float}
\usepackage{graphicx}
\usepackage{subcaption}
\usepackage{enumitem}
\usepackage{mathtools}
\usepackage{amssymb}
\usepackage{booktabs}
\usepackage[unicode=true, pdftex, colorlinks=true, linkcolor=black, citecolor=black, urlcolor=black]{hyperref}
\usepackage{listings}
\usepackage{xcolor}

\definecolor{greencomments}{rgb}{0.125,0.5,0.125}

\renewcommand{\lstlistingname}{Листинг}

\lstdefinestyle{gost}{
	language=Python,
	basicstyle=\small\ttfamily\linespread{1}\selectfont,
	keywordstyle=\color{gray!30!black}\bfseries,
	commentstyle=\color{greencomments},
	stringstyle=\color{gray!30!black},
	numberstyle=\small\color{gray!30!black},
	numbers=left,
	numbersep=12pt,
	columns=fullflexible,
	tabsize=4,
	showspaces=false,
	showstringspaces=false,
	breaklines=true,
	backgroundcolor=\color{white},
	frame=single,
	rulecolor=\color{gray!30!black},
	captionpos=t,
	escapeinside={/*@}{@*/},
	inputencoding=utf8,
	extendedchars=true,
	keepspaces=true,
	literate={а}{{\selectfont\char224}}1
	{б}{{\selectfont\char225}}1
	{в}{{\selectfont\char226}}1
	{г}{{\selectfont\char227}}1
	{д}{{\selectfont\char228}}1
	{е}{{\selectfont\char229}}1
	{ё}{{\"e}}1
	{ж}{{\selectfont\char230}}1
	{з}{{\selectfont\char231}}1
	{и}{{\selectfont\char232}}1
	{й}{{\selectfont\char233}}1
	{к}{{\selectfont\char234}}1
	{л}{{\selectfont\char235}}1
	{м}{{\selectfont\char236}}1
	{н}{{\selectfont\char237}}1
	{о}{{\selectfont\char238}}1
	{п}{{\selectfont\char239}}1
	{р}{{\selectfont\char240}}1
	{с}{{\selectfont\char241}}1
	{т}{{\selectfont\char242}}1
	{у}{{\selectfont\char243}}1
	{ф}{{\selectfont\char244}}1
	{х}{{\selectfont\char245}}1
	{ц}{{\selectfont\char246}}1
	{ч}{{\selectfont\char247}}1
	{ш}{{\selectfont\char248}}1
	{щ}{{\selectfont\char249}}1
	{ъ}{{\selectfont\char250}}1
	{ы}{{\selectfont\char251}}1
	{ь}{{\selectfont\char252}}1
	{э}{{\selectfont\char253}}1
	{ю}{{\selectfont\char254}}1
	{я}{{\selectfont\char255}}1
	{А}{{\selectfont\char192}}1
	{Б}{{\selectfont\char193}}1
	{В}{{\selectfont\char194}}1
	{Г}{{\selectfont\char195}}1
	{Д}{{\selectfont\char196}}1
	{Е}{{\selectfont\char197}}1
	{Ё}{{\"E}}1
	{Ж}{{\selectfont\char198}}1
	{З}{{\selectfont\char199}}1
	{И}{{\selectfont\char200}}1
	{Й}{{\selectfont\char201}}1
	{К}{{\selectfont\char202}}1
	{Л}{{\selectfont\char203}}1
	{М}{{\selectfont\char204}}1
	{Н}{{\selectfont\char205}}1
	{О}{{\selectfont\char206}}1
	{П}{{\selectfont\char207}}1
	{Р}{{\selectfont\char208}}1
	{С}{{\selectfont\char209}}1
	{Т}{{\selectfont\char210}}1
	{У}{{\selectfont\char211}}1
	{Ф}{{\selectfont\char212}}1
	{Х}{{\selectfont\char213}}1
	{Ц}{{\selectfont\char214}}1
	{Ч}{{\selectfont\char215}}1
	{Ш}{{\selectfont\char216}}1
	{Щ}{{\selectfont\char217}}1
	{Ъ}{{\selectfont\char218}}1
	{Ы}{{\selectfont\char219}}1
	{Ь}{{\selectfont\char220}}1
	{Э}{{\selectfont\char221}}1
	{Ю}{{\selectfont\char222}}1
	{Я}{{\selectfont\char223}}1
}
\lstset{style=gost}

\renewcommand\thesection{\arabic{section}}

\newcommand{\figimg}[2]{%
	\IfFileExists{#2}%
	{\includegraphics[width=#1]{#2}}%
	{\fbox{\parbox[c][0.75\dimexpr#1\relax][c]{#1}{\centering\small\ttfamily\detokenize{#2}}}}%
}

\begin{document}

	% ============================================================
	% ТИТУЛЬНЫЙ ЛИСТ
	% ============================================================
	\thispagestyle{empty}
	\centering
	МИНИСТЕРСТВО НАУКИ И ВЫСШЕГО ОБРАЗОВАНИЯ РОССИЙСКОЙ ФЕДЕРАЦИИ\\
	Федеральное государственное автономное образовательное учреждение высшего образования «Санкт-Петербургский политехнический университет Петра Великого»\\[10pt]
	Институт компьютерных наук и кибербезопасности\\
	Высшая школа технологий искусственного интеллекта\\
	Направление: 02.03.01 Математика и компьютерные науки\\
	\vspace{\fill}
	{\large
		ОТЧЕТ\\
		О ВЫПОЛНЕНИИ САМОСТОЯТЕЛЬНОЙ РАБОТЫ\\
		по дисциплине «Интеллектуальный анализ и визуализация данных»\\
		[20pt]
		СРС 1.4. УСТАНОВОЧНЫЙ ПРАКТИКУМ: РАБОЧИЕ МЕСТА И ТРЕКИ\\
		СРС 2.1. КОГНИТИВНЫЕ ОСНОВЫ ВОСПРИЯТИЯ ИНФОРМАЦИИ\\[10pt]
		Трек: BI-аналитика (Qlik Sense)\\
		\vspace{\fill}
	}
	
	\begin{flushleft}
		{
			{Выполнила}\hfill \hspace{5em} Колчанова И.М.\\
			Группа \hfill \hspace{5em} 5130201/40102 \\[20pt]
			{Проверил}\hfill \hspace{5em} Большаков А.А.\\
		}
	\end{flushleft}
	
	\begin{flushright}
		{<<\hrf{2em}>> \hrf{6em} 2026~г.}
	\end{flushright}
	
	\centering
	Санкт-Петербург, 2026
	
	\newpage
	
	% ============================================================
	% СОДЕРЖАНИЕ
	% ============================================================
	\justifying
	\renewcommand\contentsname{\vspace{-3.2cm}\Large СОДЕРЖАНИЕ}
	\tableofcontents
	\thispagestyle{empty}
	\newpage
	
	% ============================================================
	% ВВЕДЕНИЕ
	% ============================================================
	\section*{ВВЕДЕНИЕ}\addcontentsline{toc}{section}{ВВЕДЕНИЕ}
	
	Цель отчёта --- представить результаты самостоятельной работы по дисциплине «Интеллектуальный анализ и визуализация данных», включая выбор аналитического инструмента, рассмотрение особенностей восприятия визуальной информации и подготовку материалов для проектирования дашборда. Работа выполняется в рамках сквозного кейса --- анализа инцидентов SOC на основе учебного датасета \texttt{iavd\_incidents\_raw\_target.parquet}.
	
	В качестве профильного трека выбран трек «BI-аналитика» с использованием Qlik Sense Desktop. В отчёте зафиксирован выбор аналитического инструмента, представлены диаграммы по данным SOC, рассмотрены визуальные кодировки, предложена DIKW-проекция кейса и приведён эскиз обзора карточек инцидентов. Переработка диаграммы в формате «до --- после» и анализ траектории взгляда будут выполнены отдельно.

    \newpage
	% ============================================================
	% СРС 1.4
	% ============================================================
	\section{СРС 1.4. Установочный практикум: рабочие места и треки}
	
	\subsection{Цель и форматы}
	
	Цель работы --- выбрать профильное направление и подготовить программную среду для последующего анализа и визуализации данных.
	
	\subsection{Выбранный трек}
	
	Выбран трек \textbf{BI-аналитика} (вариант~А по умолчанию). Инструмент: Qlik Sense Desktop. Обоснование: трек ориентирован на дескриптивную аналитику и разведочный анализ данных (EDA) без написания кода, что соответствует задачам анализа исторических данных SOC и построения интерактивных дашбордов.
	
	\subsection{Чек-лист установки ПО}
	
	\subsubsection{Требования к рабочему месту}
	
	\begin{table}[H]
		\centering
		\caption{Требования к рабочему месту для трека BI-аналитика}
		\begin{tabular}{@{}ll@{}}
			\toprule
			Требование & Значение \\
			\midrule
			Процессор & 4+ ядра, x64 \\
			ОЗУ & 8 ГБ \\
			Диск & 5 ГБ для установки; SSD с запасом 20+ ГБ \\
			ОС и платформа & Windows 10/11; .NET Framework 4.8+ \\
			\bottomrule
		\end{tabular}
	\end{table}
	
	
	\subsection{Форма фиксации трека}
	
	\begin{table}[H]
		\centering
		\caption{Форма фиксации трека}
		\begin{tabular}{@{}ll@{}}
			\toprule
			Параметр & Значение \\
			\midrule
			Номер группы & 5130201/40102 \\
			Состав (ФИО) &  \\
			Трек & BI-аналитика (Qlik Sense) \\
			\bottomrule
		\end{tabular}
	\end{table}
	
	\subsection{Критерии завершения занятия}
	
	Для дальнейшей работы необходимы доступ к выбранному инструменту BI-аналитики, возможность загрузки набора данных и проверка создания тестовой визуализации. В таблицах выше зафиксированы требования к программной среде и выбранное направление.
	
	\subsection{Ответы на контрольные вопросы}
	
	\begin{enumerate}[label=\arabic*.]
		\item Версии библиотек в окружении должны совпадать с версиями из \texttt{pyproject.toml}, потому что предобученная модель курса сохранена в формате \texttt{joblib}, и при несовпадении версий \texttt{scikit-learn}, \texttt{numpy} и других зависимостей десериализация модели завершится ошибкой или даст некорректные предсказания.
		\item Если шаг чек-листа не выполняется, необходимо сообщить преподавателю до 12.09.2026 (занятие~1.5). Переносить это на время практических занятий нельзя, потому что практикум 1.4 ограничен по времени (50 минут), а нерешённые проблемы с окружением блокируют выполнение последующих работ.
		\item Фиксация трека на занятии~1.4 даёт команде определённость: выбран инструмент, распределены роли, сформировано понимание дальнейших шагов. Это позволяет приступить к работе с датасетом без организационных задержек.
	\end{enumerate}
	
	% ============================================================
	% СРС 2.1
	% ============================================================
	\section{СРС 2.1. Когнитивные основы восприятия информации}
	
	\subsection{Цель работы}
	
	Цель --- применить принципы когнитивной психологии восприятия к проектированию обзора карточек инцидентов SOC: выбрать ключевые поля, обосновать кодировки по иерархии Кливленда--МакГилла, построить эскиз обзора.
	
	\subsection{Построенные диаграммы}
	
	В Qlik Sense Desktop построены три диаграммы на основе датасета \texttt{iavd\_incidents\_raw\_target.parquet}.
	
	\subsubsection{Диаграмма 1. Доля ложных срабатываний по уровням критичности}
	
	\textbf{Измерение:} \texttt{Критичность}. \textbf{Мера:} \texttt{Count({<Вердикт=\{'False Positive'\}>} EntityId) / Count(EntityId)}. \textbf{Тип:} линейчатая диаграмма (Bar chart).
	
	\begin{figure}[H]
		\centering
		\figimg{0.9\textwidth}{Pictures/chart1_criticality.png}
		\caption{Доля ложных срабатываний внутри каждого уровня критичности}
		\label{fig:chart1}
	\end{figure}
	
	\textbf{Наблюдение.} Наибольшая доля ложных срабатываний наблюдается в категориях «Отсутствует» (около 74\,\%) и «Средняя» (около 71,2\,\%), наименьшая --- в категории «Критическая» (около 25\,\%). При этом зависимость доли False Positive от критичности не является строго монотонной. Показатель описывает состав карточек внутри каждой категории, но сам по себе не определяет вероятность реальной атаки.
	
	\subsubsection{Диаграмма 2. Распределение карточек по вердикту аналитика}
	
	\textbf{Измерение:} \texttt{Вердикт}. \textbf{Мера:} относительная частота карточек по категориям вердикта. \textbf{Тип:} линейчатая диаграмма (Bar chart). Значения представлены долями от общего числа карточек.
	
	\begin{figure}[H]
		\centering
		\figimg{0.9\textwidth}{Pictures/chart2_verdict.png}
		\caption{Распределение карточек по вердикту аналитика, доля от общего числа}
		\label{fig:chart2}
	\end{figure}
	
	\textbf{Наблюдение.} Почти две трети карточек (около 62,7\,\%) имеют вердикт False Positive. Категория «Не указан» составляет около 22,8\,\%, «Решён» --- 11,9\,\%, «Дубликат» --- 2,4\,\%; присутствуют и другие малочисленные категории. Значительная доля ложных срабатываний указывает на потенциально высокую нагрузку на аналитиков и мотивирует дальнейшую проверку способов приоритизации.
	
	\subsubsection{Диаграмма 3. Динамика создания алертов по часам суток}
	
	\textbf{Измерение:} \texttt{Hour(Создан)}. \textbf{Мера:} \texttt{Count(EntityId)}. \textbf{Тип:} линейчатая диаграмма (Bar chart).
	
	\begin{figure}[H]
		\centering
		\figimg{0.9\textwidth}{Pictures/chart3_hours.png}
		\caption{Динамика создания алертов по часам суток, шт.}
		\label{fig:chart3}
	\end{figure}
	
	\textbf{Наблюдение.} Пик создания карточек приходится на 8:00 (19\,118 карточек), минимум --- на 0:00 (3\,591 карточка). Ночные часы (0--6) дают 19,1\,\% карточек. Это указывает на наличие регистрируемых событий в течение суток, но не позволяет само по себе делать вывод о фактическом графике работы сотрудников SOC.
	
	\subsection{Обоснование выбора кодировок}
	
	Для сравнения \textbf{долей ложных срабатываний} внутри уровней критичности используется длина столбцов от общей нулевой линии. Такое кодирование позволяет сопоставлять значения без необходимости сравнивать площади или углы. Уровни критичности представлены подписями категорий; в текущей диаграмме используется единый цвет столбцов, поэтому цвет не кодирует степень критичности. Для будущего интерфейса цвет можно использовать как дополнительный сигнал, не заменяя текстовые подписи.
	
	В проектируемом обзоре карточек для \textbf{приоритизации} предлагается сортировка по убыванию оценки риска и визуальное выделение карточек, требующих первоочередного внимания. Вердикт аналитика при этом необходимо отображать отдельной текстовой меткой: положение в очереди не является кодировкой самого вердикта. Насыщенный сигнальный цвет целесообразно применять избирательно, сохраняя нейтральное оформление остальных элементов.
	
	Группировка элементов выполняется по гештальт-паттернам: KPI-показатели размещаются на общей плашке (паттерн «общая область» + «близость»), а связанные поля \texttt{src} и \texttt{dst} --- вплотную друг к другу (паттерн «близость»), что соответствует принципу «рядом = одна пара события».
	
		\subsection{DIKW-проекция кейса}
	
	Модель DIKW (Data -- Information -- Knowledge -- Decision) применяется к сквозному кейсу анализа инцидентов SOC. Переход от исходных записей к интерпретации и решению требует контекста и проверки предположений.
	
	\begin{table}[H]
		\centering
		\caption{DIKW-проекция кейса SOC}
		\begin{tabular}{@{}p{2.5cm}p{11cm}@{}}
			\toprule
			Уровень & Содержание \\
			\midrule
			\textbf{Data} & Строки логов SIEM: \texttt{timestamp}, \texttt{CMDLINE}, \texttt{processName}, \texttt{src}, \texttt{dst}, \texttt{suser}. Пример: \texttt{2026-08-14 02:31:07}, \texttt{powershell.exe -nop -w hidden -enc JABzAD0A...}, \texttt{10.12.4.87}, \texttt{203.0.113.147}, \texttt{CORP\textbackslash a.ivanov}. \\
			\textbf{Information} & Обогащённая карточка инцидента INC-2026-004817: актив --- рабочая станция 10.12.4.87; пользователь --- \texttt{CORP\textbackslash a.ivanov}; время --- 02:31 (ночь); правило корреляции SIEM --- «Подозрительное выполнение PowerShell с обфускацией»; критичность --- высокая; \texttt{objectValue} --- ключ автозагрузки \texttt{HKCU\textbackslash...\textbackslash Run\textbackslash OneDriveStandaloneUpdater}. \\
			\textbf{Knowledge} & Проверенная закономерность: сочетание «ночное время + скрытый запуск PowerShell (\texttt{-w hidden -enc}) + ключ автозагрузки + внешний \texttt{dst}» связано с повышенным риском закрепления вредоносного кода. Область применимости: учётная запись обычного сотрудника, не администратора. Ограничение: возможны ложные срабатывания при регламентных работах. \\
			\textbf{Decision} & Эскалировать карточку на L2, запросить телеметрию по хосту 10.12.4.87, проверить историю активности \texttt{a.ivanov}, изолировать хост при подтверждении. \\
			\bottomrule
		\end{tabular}
	\end{table}
	
	\subsubsection{Какой контекст отсутствует}
	
	Карточка инцидента отражает ограниченный набор сведений об одном событии и не содержит полного контекста. Для обоснованного вердикта не хватает:
	
	\begin{itemize}
		\item истории активности пользователя \texttt{a.ivanov} (было ли это раньше?);
		\item информации о регламентных работах в 02:31;
		\item связанных событий до и после инцидента;
		\item привилегий учётной записи (администратор или обычный сотрудник?);
		\item внутренний или внешний \texttt{dst} (203.0.113.147 --- учебный внешний адрес).
	\end{itemize}
	
	\subsubsection{Цена FP и FN}
	
	В рассматриваемом кейсе пропуск настоящего инцидента (FN) может иметь более серьёзные последствия, чем отдельная ложная тревога (FP).
	
	\begin{table}[H]
		\centering
		\caption{Цена ошибок модели в кейсе SOC}
		\begin{tabular}{@{}p{2.5cm}p{4.5cm}p{6.5cm}@{}}
			\toprule
			Ошибка & Что значит & Последствия \\
			\midrule
			\textbf{FP} & Модель сказала «инцидент», а его нет & Дополнительная проверка (5--10 минут), задержка других карточек, «усталость от алертов», риск ошибочной блокировки. При условных 1\,274 карточках за смену и доле FP 62,7\,\% это составило бы примерно 799 карточек с таким вердиктом; данный расчёт иллюстративен. \\
			\textbf{FN} & Модель сказала «легитимно», а это инцидент & Задержка обнаружения, развитие атаки, компрометация данных, простой бизнеса, репутационные потери, рост ущерба. \\
			\bottomrule
		\end{tabular}
	\end{table}
	
	Следовательно, при выборе порога принятия решения необходимо учитывать сравнительную стоимость FN и FP. Снижение риска пропуска атак важно, однако чрезмерное увеличение числа FP также создаёт нагрузку на аналитиков; оптимальный баланс требует проверки на данных.
	
	\subsection{Эскиз обзора карточек}
	
	Эскиз обзора карточек инцидентов представлен на рисунке~\ref{fig:sketch}. Он содержит четыре смысловые зоны: KPI-блок, очередь карточек с подсветкой, плашку «Требует вердикта аналитика» и таймлайн алертов.
	
	\begin{figure}[H]
		\centering
		\figimg{1\textwidth}{Pictures/sketch.pdf}
		\caption{Эскиз обзора карточек инцидентов}
		\label{fig:sketch}
	\end{figure}	
	
	\subsection{Тезисы по квартету Энскомба}
	
	\begin{enumerate}[label=\arabic*.]
		\item Квартет Энскомба показывает, что четыре набора данных с идентичными сводными статистиками (средние, дисперсии, корреляция, линия регрессии) могут иметь радикально разную структуру: линейное облако, криволинейную связь, выброс, вертикаль.
		\item Сводные статистики необходимы, но недостаточны: они скрывают форму связи, разброс, влияние отдельных точек и наличие кластеров. Для вердикта по карточке SIEM нужны её факты и контекст, а не только агрегаты.
		\item Вывод для аналитика SOC: одинаковые средние число алертов и доля критических могут скрывать разную динамику; временной график помогает её проверить.
	\end{enumerate}

    \newpage
	% ============================================================
	% ЗАКЛЮЧЕНИЕ
	% ============================================================
	\section*{ЗАКЛЮЧЕНИЕ}\addcontentsline{toc}{section}{ЗАКЛЮЧЕНИЕ}
	
	В работе зафиксирован выбор BI-трека и рассмотрены требования к рабочему месту. На примере учебного набора данных SOC проанализированы доли ложных срабатываний по критичности, распределение вердиктов и число карточек по часам суток. Сформулированы принципы визуального кодирования, представлена DIKW-проекция кейса, рассмотрена цена ошибочных решений и подготовлен эскиз интерфейса обзора инцидентов.

	Следующими этапами являются переработка одной из диаграмм по принципу «до --- после» и проверка логики восприятия интерфейса с помощью траектории взгляда. Эти результаты пока не включены в отчёт как выполненные задания.

	% ============================================================
	% СПИСОК ЛИТЕРАТУРЫ
	% ============================================================
	\renewcommand{\bibname}{\vspace{-2cm}\large СПИСОК ИСПОЛЬЗОВАННЫХ ИСТОЧНИКОВ}
	\addcontentsline{toc}{section}{СПИСОК ИСПОЛЬЗОВАННЫХ ИСТОЧНИКОВ}
	
	\begin{thebibliography}{9}
		\bibitem{tema}
		Большаков А.А. Интеллектуальный анализ и визуализация данных: методические указания к занятию~1.4. СПб.: СПбПУ, 2026.
		
		\bibitem{wolfram}
		Большаков А.А. Когнитивные основы восприятия информации: презентация к занятию~2.1. СПб.: СПбПУ, 2026.
		
		\bibitem{habr}
		Нафлик К. Данные: визуализируй, расскажи, используй. Сторителлинг в аналитике. М.: Эксмо, 2025.
		
		\bibitem{cleveland}
		Cleveland W.S., McGill R. Graphical Perception: Theory, Experimentation, and Application to the Development of Graphical Methods // Journal of the American Statistical Association. 1984. Vol.~79, No.~387. P.~531--554.
		
		\bibitem{cowan}
		Cowan N. The magical number 4 in short-term memory: A reconsideration of mental storage capacity // Behavioral and Brain Sciences. 2001. Vol.~24, No.~1. P.~87--114.
	\end{thebibliography}
	
\end{document}
